\begin{frame}{결과 읽기: 문맥이 비슷한 단어가 가깝다}
  \begin{itemize}
    \item ``동물'' 관련(동물이다, 사자는, 강아지는)이 ``과일'' 관련
      (과일이다, 포도는, 바나나는)보다 \textbf{대체로} 코사인 유사도가
      높음.
    \item ``고양이는''이 실제로 등장한 문맥(동물 문장)과 비슷한
      문맥에 나온 단어일수록 \textbf{벡터 공간에서도 가까움}.
    \item 말뭉치가 아주 작아 완벽하진 않다(예: ``사과는''가 예상보다
      높게) --- 데이터가 수십억 단어로 커지면 노이즈가 줄고,
      ``의미적 유사도 = 벡터의 각도''가 선명해진다.
  \end{itemize}
  \vskip0.6em
  \begin{block}{왜 코사인 유사도인가 (14.1 PCA와 같은 이유)}
    벡터의 \textbf{``방향''이 의미}이고, \textbf{``길이''는 (주로)
    빈도}라서, 비슷한 단어를 찾을 때 각도(= 코사인)를 보고 길이는
    무시하는 것이 맞다.
  \end{block}
\end{frame}
